
\section{Bialgebras}\label{sec:bialgebras}

Here we give the $2$-categorical definition of an object of bialgebras, extending Ulmer's definition \cite{Ulmer1977Bialgebras} from $\CAT$ to any $2$-category with PIE limits. The choice of structural maps of a bialgebraic object is captured by inserters and equations are enforced by equifiers. We then develop a multi-sorted, $2$-categorical syntactic framework (signature pairs and bialgebraic theories) which produces the data of this definition automatically from the syntax itself.

\subsection{Bialgebras in a \texorpdfstring{$2$}{2}-Category}

\begin{definition}[Bialgebraic Data and Their Bialgebras]\label{def:general_bialgebras}
   Let $\K$ be a $2$-category with inserters and equifiers, and let $c$ be an object in $\K$. A \textbf{bialgebraic datum} on $c$ is a tuple
   $$
   \mathbb{B} = \big(c,\ f,g\colon c\to x,\ \ f',g'\colon c\to y,\ \ \alpha,\beta\colon f'u'\Rightarrow g'u'\big),
   $$
   where $(u'\colon \Ins(f,g)\to c,\ \phi\colon fu'\Rightarrow gu')$ is the inserter of $f$ and $g$. The morphisms $f,g,f',g'$ are called the \textbf{supports} of $\mathbb{B}$:
   \begin{itemize}
       \item $f$ is the \textbf{function arity support} and $g$ the \textbf{function coarity support};
       \item $f'$ is the \textbf{equational arity support} and $g'$ the \textbf{equational coarity support}.
   \end{itemize}
   Thus $f,f'$ are the \textbf{arity supports}, $g,g'$ the \textbf{coarity supports}, $f,g$ the \textbf{function supports} and $f',g'$ the \textbf{equational supports}. The pair $(\alpha,\beta)$ is called the \textbf{equations} of $\mathbb{B}$.
   \begin{enumerate}
       \item The \textbf{object of pre-bialgebras} of $\mathbb{B}$ is the inserter $u'\colon \Ins(f,g)\to c$.
       \item The \textbf{object of bialgebras} of $\mathbb{B}$ is the equifier
       $$
       \mathrm{Bialg}(\mathbb{B})\coloneqq \Eq(\alpha,\beta)\xrightarrow{e}\Ins(f,g).
       $$
       \item By $u_\B = u\colon\BialgB\to c$ we denote the composite $u\coloneqq u'e$.
   \end{enumerate}
   We use capital letters $(\C, F,G,\ F',G',\ \alpha,\beta)$ to indicate the concrete nature of $\K$, for instance, when $\K = \CAT$ or $\K =\V$-$\CAT$. 
\end{definition}

Originally, Ulmer considered families of pairs of support morphisms $(f_\mu,g_\mu\colon c\to x_\mu)_{\mu}$ and equations $(\alpha_r,\beta_r)_{r}$ indexed by classes. When the classes are sets, this data can be turned into one pair by taking products, $f=(f_\mu)_\mu, g= (g_\mu)_\mu\colon c\to\prod_\mu x_\mu$, and similarly for the equations. For $\K=\CAT$ this recovers Ulmer's definition \cite{Ulmer1977Bialgebras} when the indexing families are small. We sidestep the size issues via the Grothendieck universe approach.

The syntax devised below produces bialgebraic data automatically. We develop the notions of a signature pair $(\Sigma,\sigma)$, a $\sigma$-theory $T$ and a $\Sigma$-model $M$, which together determine a bialgebraic datum on $M_0$. The associated object of bialgebras is then $M^T$ (Definition~\ref{def:objectOfTBialgebras}).

\begin{remark}\label{rmk:AdjointsInducingIsosBetweenBialgebras}
    Consider a bialgebraic datum $\B = (c,\ f,g,\ f',g',\ \alpha,\beta)$ in $\K$ with $g,g'$ having left adjoints $l$ and $l'$, respectively. Note that $\BialgB$ is isomorphic to $\Bialg(\B')$ over $c$, where $\B' = (c,\ c\xrightarrow{lf,\id_c} c,\ c\xrightarrow{l'f',\id_c}c,\ l'f'u''\xRightarrow{\alpha',\beta'} u''\colon \Ins(lf,\id_c)\to c)$ and $\alpha'$ and $\beta'$ are transformations induced by the adjunctions: Notice first that the adjunction $l\dashv g$ induces an isomorphism $\Ins(f,g)\cong \Ins(lf,\id_c)$ over $c$. This allows one to consider $\alpha'',\beta''\colon f'u''\Rightarrow g'u''$, where $u''\colon \Ins(lf,\id_c)\to c$ is the inserter, and whose mates we denote by $\alpha',\beta'\colon l'f'u''\Rightarrow u''$.

    The previous construction also works in the case that $g$ and $g'$ only have $f$- and $f'$-relative left adjoints $r,r'\colon c\to c$, respectively.\footnote{If $F\colon I\to \D$, $J\colon I\to \C$ and $G\colon \D\to \C$ are functors, then we say that $F$ is a $J$-relative left adjoint of $G$ if there is an isomorphism $\D(Fx,y)\cong \C(Jx,Gy)$ natural in $x$ and $y$. One can then define when $f$ is a $j$-relative left adjoint of $g$ in a $2$-category $\K$ via representable $2$-functors.}
\end{remark}

\subsection{Syntax: Signature Pairs and Theories}
To motivate the syntactic constructions in this section, consider the syntax needed to formulate internal monoids in a lax monoidal setting:
\[
\Sigma = (c, \quad cc\xrightarrow{\square} c,\ ()\xrightarrow{I} c, \quad (x\square y)\square z\xrightarrow{\alpha^{xyz}}x\square (y\square z),\ x\square I\xrightarrow{\rho^x} x,\ I\square x\xrightarrow{\lambda^x} x),
\]
a relative $1$-signature describing the monoid's multiplication and unit operations
\[
\sigma = (s\colon c, \quad \mu\colon s\square s\to s, \quad u\colon I\to s),
\]
and the associativity and unit equations $T$. This information is packaged within a signature pair $(\Sigma,\sigma)$ and a $\sigma$-theory $T$.

Given a function $f \colon X \to Y$ between sets, we say that the set $X$ is \emph{typed} in $Y$ and write $x \colon f(x)$ for each element $x$ in $X$.

\begin{definition}[The $2$-Signature $\Sigma$]
   A \textbf{$2$-signature} $\Sigma = (C,\mathcal{F}, \mathcal{T})$ consists of the following data:
   \begin{enumerate}
       \item A small set $C$ of \textbf{category symbols}, often denoted by $c,d,c_1,c_2,\dots$.
       \item A typed small set $\mathcal{F} \to C^* \times C$ of \textbf{functor symbols} $B \colon c_1 \cdots c_n \to d$, where $C^*$ is the free monoid over $C$.
       \item A typed set $V \to C$, with countably infinite fibers, of \textbf{variable symbols} $x,y,z,\dots$. From these variables, we recursively construct a typed set $\Term^\Sigma_V \to C$ of \textbf{$\Sigma$-terms}.\footnote{A term $t\in \Term_V^\Sigma$ is either a variable $x\colon c$ in $V$, a constant $B\colon ()\to c$ (functor symbol with empty domain) or $t = B(t_0,\ldots, t_n)\colon c$ for some previously constructed terms $t_0\colon c_0,\ldots, t_n\colon c_n$ and a functor symbol $B\colon c_0\cdots c_n\to c\in \mathcal{F}$, where $n\in \N$ (with $0\in \N$).} The set of \textbf{contexts} $\mathrm{Cont}_V \subset V^*$ consists of sequences where no variable occurs more than once. A term $t \in \Term_V^\Sigma$ is in a context $\bar{x} \in \mathrm{Cont}_V$ if all variables appearing in $t$ occur in the sequence $\bar{x}$.
       \item A small set $\mathcal{T} \to \mathrm{Cont}_V \times \Term^\Sigma_V \times \Term^\Sigma_V$ of \textbf{transformation symbols} $\alpha^{\bar{x}} \colon t_1 \to t_2$, where $t_1, t_2$ are $\Sigma$-terms of the same type in context $\bar{x}$.
   \end{enumerate}
\end{definition}
\begin{example}\label{ex:2catInducesSignature}
Each small $2$-category $\A$ induces a $2$-signature $\Sigma_\A = (C, \mathcal{F}, \mathcal{T})$ as follows:
\begin{itemize}
    \item \textbf{Category symbols:} The set $C \coloneqq \Obj(\A)$. For each $a \in \Obj(\A)$, we fix a variable symbol $x_a$ of type $a$.
    \item \textbf{Functor symbols:} Each $1$-cell $f \colon a \to b$ in $\A$ gives a functor symbol $f \colon a \to b$.
    \item \textbf{Transformation symbols:} We include three families of transformation symbols:
    \begin{itemize}
        \item A transformation symbol $\eta^{x_a} \colon f(x_a) \to g(x_a)$ for each $2$-cell $\eta \colon f \Rightarrow g \colon a \to b$ in $\A$.
        \item A transformation symbol $\alpha^{x_a}_{g,f} \colon (gf)(x_a) \to g(f(x_a))$ for each composable pair $a \xrightarrow{f} b \xrightarrow{g} c$ in $\A$.
        \item A transformation symbol $\beta^{x_a}_a \colon \id_a(x_a) \to x_a$ for each object $a \in \Obj(\A)$.
    \end{itemize}
\end{itemize}
\end{example}

We define the internal algebraic structures on top of this background categorical universe by introducing a relative $1$-signature.

\begin{definition}[The $1$-Signature $\sigma$]
   A pair $(\Sigma,\sigma)$ is called a \textbf{signature pair} if $\Sigma = (C,\mathcal{F},\mathcal{T})$ is a $2$-signature and the \textbf{$1$-signature} $\sigma = (S, \Phi)$ relative to $\Sigma$ consists of:
   \begin{enumerate}
       \item A typed small set of \textbf{sorts} $S \to C$, assigning each sort $s \in S$ a type $c_s \in C$.
       \item A typed small set of \textbf{function symbols} $\Phi \to \Term_{S}^\Sigma \times \Term_{S}^\Sigma$. Each $f \in \Phi$ is associated with a tuple $(t_1,t_2)$, denoted $f\colon t_1 \to t_2 \colon c_f$, where $t_1, t_2\in \Term_S^\Sigma$ are terms of type $c_f$.
   \end{enumerate}
\end{definition}

We recursively generate a graph of terms in order to formulate equations between these operations.

\begin{definition}[The Term Graph]
   Let $(\Sigma,\sigma)$ be a signature pair. We construct the \textbf{term graph} $G_\sigma$ in the following recursive way:
   \begin{itemize}
       \item The set of vertices of $G_\sigma$ is the set $\Term^\Sigma_S$ of $\Sigma$-terms over the set of sorts $S$ of $\sigma$.
       \item An arrow of $G_\sigma$ is recursively defined by:
       \begin{itemize}
           \item Each function symbol $f \colon t_1 \to t_2$ of $\sigma$ forms an arrow $t_1 \to t_2$ in $G_\sigma$.
           \item For each transformation symbol $\alpha^{\bar{x}} \colon t_1 \to t_2$ in $\Sigma$, where $\bar{x} = x_1\cdots x_n$, we set $\alpha_{t_1',\dots, t_n'} \colon t_1^{\bar{x}}(t_1',\dots, t_n') \to t_2^{\bar{x}}(t_1',\dots, t_n')$ as an arrow in $G_\sigma$ for each sequence of terms $t_i' \in \Term_S^\Sigma$ of type $c_{x_i}$.\footnote{The term $t_i^{\bar{x}}(t_1',\ldots, t_n')$ refers to recursively defined substitution, where $t_j'$ is substituted for $x_j$ in $t_i$ for $i = 1,2$ and $j\leq n$.}
           \item For any functor symbol $B \colon c_0\cdots c_n \to d$ in $\Sigma$, a sequence of terms $t_j \colon c_j$, and an arrow $g \colon t_i \to t_i'$ in $G_\sigma$, the symbol sequence
           \[
           B(t_0,\dots, t_{i-1}, g, t_{i+1}, \dots, t_n) \colon B(t_0,\dots, t_i, \dots, t_n) \to B(t_0,\dots, t_i', \dots, t_n)
           \]
           defines an arrow in $G_\sigma$.\footnote{We use the indexing starting from $0$ in terms $B(t_0,\ldots, t_n)$ to indicate that $B$ is not a constant symbol. If the indexing starts from $1$, it implies the possibility that $B(t_1,\ldots, t_n)$ is a constant symbol with the choice $n = 0$.}
       \end{itemize}
   \end{itemize}
\end{definition}

Using this term graph, we define equations and classify theories based on their arity and coarity behaviors.

\begin{definition}[Arity, Coarity, and Bialgebraic Theories]
   Let $(\Sigma,\sigma)$ be a signature pair. A parallel path $(p, q \colon t_1 \to t_2\colon c)$ in the term graph $G_\sigma$ is called a \textbf{$\sigma$-equation}, and a set $T$ of $\sigma$-equations is called a \textbf{bialgebraic $\sigma$-theory}. We define the arity and coarity terms associated with $T$ as follows:
   \begin{enumerate}
       \item \textbf{Function Arity and Coarity:} For each function symbol $f \colon t_1 \to t_2$ in the $1$-signature $\sigma$, the domain term $t_1$ is a \textbf{function arity term} and the codomain term $t_2$ is a \textbf{function coarity term}.
       \item \textbf{Equational Arity and Coarity:} For each equation $(p, q \colon t_1' \to t_2')$ in the theory $T$, the domain term $t_1'$ is an \textbf{equational arity term} and the codomain term $t_2'$ is an \textbf{equational coarity term}.
       \item \textbf{Arity Terms:} The collection of all function arity terms and equational arity terms is the set of \textbf{arity terms}.
       \item \textbf{Coarity Terms:} The collection of all function coarity terms and equational coarity terms is the set of \textbf{coarity terms}.
   \end{enumerate}
   A bialgebraic $\sigma$-theory $T$ is called \textbf{algebraic} if all coarity terms consist solely of sorts $s \in S$. Similarly, the theory $T$ is called \textbf{coalgebraic} if all arity terms are sorts.
\end{definition}

\subsection{Semantics: Models and the Object of Bialgebras}

We now define the semantics of these signatures by interpreting them as models, operations, and equations within an arbitrary PIE-complete $2$-category. We fix a $2$-category $\K$ with finite products and a signature pair $(\Sigma, \sigma) = ((C,\mathcal{F}, \mathcal{T}),(S,\Phi))$ with a $\sigma$-theory $T$. 

\begin{definition}[$\Sigma$-Models]
   A \textbf{$\Sigma$-model} $M$ in $\K$ assigns:
   \begin{enumerate}
       \item An object $M_c$ in $\K$ to each category symbol $c \in C$. For any typed sequence $\bar{x} = (x_1 \colon c_1, \dots, x_n \colon c_n)$ with $c_1,\ldots, c_n\in C$, we define the product object $M_{\bar{x}} = \prod_{i=1}^n M_{c_i}$ in $\K$.
       \item A $1$-cell $M(B) \colon M_{\bar{c}} \to M_d$ in $\K$ for each functor symbol $B \colon \bar{c} \to d \in \mathcal{F}$. For any $\Sigma$-term $t \colon c$ in a context $\bar{x}$, we recursively construct the term $1$-cell $M_{\bar{x}}(t) \colon M_{\bar{x}} \to M_c$.\footnote{Consider $t\in \Term^\Sigma_V$ in context $\bar{x}$. If $t = x_i$, then $M_{\bar{x}}(t) = \pi_{i}\colon M_{\bar{x}}\to M_{c_i}$ and if $t = B(t_1,\ldots, t_n)$, then $M_{\bar{x}}(t) = M(B)(M_{\bar{x}}(t_1),\ldots, M_{\bar{x}}(t_n))$.}
       \item A $2$-cell $M(\alpha) \colon M_{\bar{x}}(t_1) \Rightarrow M_{\bar{x}}(t_2)$ in $\K$ for each transformation symbol $\alpha^{\bar{x}} \colon t_1 \to t_2 \in \mathcal{T}$.
   \end{enumerate}
\end{definition}
\begin{example}\label{ex:2FunctorInucingAModel}
Let $\A$ and $\K$ be $2$-categories with $\A$ small, and let $\Sigma_\A$ be the $2$-signature of Example~\ref{ex:2catInducesSignature}. Each strict $2$-functor $F \colon \A \to \K$ induces a $\Sigma_\A$-model $M_F$ in $\K$, defined by:
\[
M_F(a) = Fa, \quad M_F(f) = Ff, \quad M_F(\eta^{x_a}) = F\eta, \quad M_F(\alpha_{g,f}^{x_a}) = \id_{F(gf)}, \quad M_F(\beta^{x_a}_a) = \id_{\id_{Fa}}.
\]
\end{example}

Between different $\Sigma$-models, we can define three natural notions of morphism.

\begin{definition}[Lax, Oplax, and Strong Morphisms]\label{def:laxStrongMorphisms}
   Let $M$ and $N$ be $\Sigma$-models in $\K$. A \textbf{lax morphism} $(F,\eta) \colon M \to N$ consists of the following data:
   \begin{enumerate}
       \item A $1$-cell $F_c \colon M_c \to N_c$ for each category symbol $c$ in $\Sigma$. We define $F_{\bar{x}} = \prod_{i\leq n} F_{c_i}\colon M_{\bar{x}}\to N_{\bar{x}}$ for every typed sequence $\bar{x} = (x_1\colon c_1,\ldots, x_n\colon c_n)$, where $c_1,\ldots, c_n$ are category symbols of $\Sigma$.
       \item A $2$-cell $\eta^B \colon N(B) \circ F_{\bar{c}} \Rightarrow F_d \circ M(B)$ for every functor symbol $B \colon \bar{c} \to d$ in $\Sigma$. We recursively extend this to define a $2$-cell $\eta_{t, \bar{x}} \colon N_{\bar{x}}(t) \circ F_{\bar{x}} \Rightarrow F_c \circ M_{\bar{x}} (t)$ for every $\Sigma$-term $t \colon c$ in a context $\bar{x} = x_1\cdots x_n$ as follows:
       {\small
   \begin{equation}\label{eq:fromPathToATransformation}
       \eta_{t,\bar{x}} = \begin{cases}
   \id\colon \pi_i\circ F_{\bar{x}}\Rightarrow F_{c_i}\circ \pi_i, & \text{if $t = x_i$ for $i\leq n$} \\
   N(B)(N_{\bar{x}}(t_1)F_{\bar{x}},\ldots, N_{\bar{x}}(t_n)F_{\bar{x}}) \\
\quad \xRightarrow{N(B)*(\eta_{t_1,\bar{x}},\dots, \eta_{t_n, \bar{x}})} N(B)F_{\bar{c}}(M_{\bar{x}}(t_1),\dots, M_{\bar{x}}(t_n))\\
   \quad \xRightarrow{\eta^B(M_{\bar{x}}(t_1),\dots, M_{\bar{x}}(t_n))} F_c M(B)(M_{\bar{x}}(t_1),\dots, M_{\bar{x}}(t_n)) & \text{if $t = B(t_1,\dots, t_n)$}
\end{cases}
\end{equation}
}
       \item For every transformation symbol $\theta^{\bar{x}} \colon t_1 \to t_2$ of type $c$ in $\Sigma$, we require that the diagram
       $$
       \begin{tikzcd}[ampersand replacement=\&]
       N_{\bar{x}}(t_1) \circ F_{\bar{x}} \arrow[r, "\eta_{t_1, \bar{x}}", Rightarrow] \arrow[d, "N(\theta)*F_{\bar{x}}"', Rightarrow] \& F_c \circ M_{\bar{x}}(t_1) \arrow[d, "F_c*M(\theta)", Rightarrow] \\
       N_{\bar{x}}(t_2) \circ F_{\bar{x}} \arrow[r, "\eta_{t_2, \bar{x}}"' , Rightarrow]                                               \& F_c \circ M_{\bar{x}}(t_2)
       \end{tikzcd}
       $$
       commutes.
   \end{enumerate}
   Dually, we define an \textbf{oplax morphism} by reversing the direction of the $2$-cells $\eta^B$. We say that $(F,\eta)$ is a \textbf{strong morphism} if it is a lax morphism where each $2$-cell $\eta^B$ is an isomorphism. A lax morphism $(F,\eta)\colon M\to N$ is called \textbf{strict} if $\eta^B$ is an identity for each $\Sigma$-functor symbol $B$.
\end{definition}

\begin{proposition}
   Let $(F,\eta) \colon M \to N$ and $(G,\theta) \colon N \to P$ be lax morphisms of $\Sigma$-models in $\K$. Then the composite $G \circ F$ is a lax morphism. Furthermore, oplax and strong morphisms are closed under composition.
\end{proposition}
\begin{proof}
   We sketch the case for lax morphisms. We define the components $(G \circ F)_c = G_c \circ F_c$ and define the structural $2$-cell $\phi^B$ for each functor symbol $B \colon \bar{c} \to d$ as the composite:
   $$
   \phi^B \colon P(B) \circ (GF)_{\bar{c}} \xRightarrow{\theta^B*F_{\bar{c}}} G_d \circ N(B) \circ F_{\bar{c}} \xRightarrow{G_d * \eta^B} (GF)_d \circ M(B)
   $$
   Via induction, for any $\Sigma$-term $t \colon c$ in a context $\bar{x}$, we obtain the factorization:
   \[
   \phi_{t, \bar{x}} \colon P_{\bar{x}}(t) \circ (GF)_{\bar{x}} \xRightarrow{\theta_{t, \bar{x}}*F_{\bar{x}}} G_c \circ N_{\bar{x}}(t) \circ F_{\bar{x}} \xRightarrow{G_c*\eta_{t, \bar{x}}} (GF)_c \circ M_{\bar{x}}(t)
   \]
   Consider a transformation symbol $\alpha^{\bar{x}} \colon t_1 \to t_2\colon c$ in $\Sigma$. Now the diagram
   \[
   \begin{tikzcd}[column sep = 2cm, ampersand replacement=\&]
   P_{\bar{x}}(t_1) \circ (GF)_{\bar{x}} \arrow[d, "P(\alpha)*(GF)_{\bar{x}}" description, Rightarrow] \arrow[rr, "\phi_{t_1, \bar{x}}", Rightarrow, bend left] \arrow[r, "\theta_{t_1, \bar{x}}*F_{\bar{x}}" description, Rightarrow] \& G_c \circ N_{\bar{x}}(t_1) \circ F_{\bar{x}} \arrow[r, "G_c*\eta_{t_1, \bar{x}}" description, Rightarrow] \arrow[d, "G_c*N(\alpha)*F_{\bar{x}}" description, Rightarrow] \& (GF)_c \circ M_{\bar{x}}(t_1) \arrow[d, "(GF)_c*M(\alpha)" description, Rightarrow] \\
   P_{\bar{x}}(t_2) \circ (GF)_{\bar{x}} \arrow[r, "\theta_{t_2, \bar{x}}*F_{\bar{x}}"' description, Rightarrow] \arrow[rr, "\phi_{t_2, \bar{x}}"', Rightarrow, bend right]                                        \& G_c \circ N_{\bar{x}}(t_2) \circ F_{\bar{x}} \arrow[r, "G_c*\eta_{t_2, \bar{x}}"' description, Rightarrow]                                                                \& (GF)_c \circ M_{\bar{x}}(t_2)
   \end{tikzcd}
   \]
   commutes. Thus, the composite $G \circ F$ is a lax morphism of $\Sigma$-models.
\end{proof}

To obtain the $2$-category of $\Sigma$-models and lax morphisms, and its sub-$2$-category of strong morphisms, we define the $2$-cells. A \textbf{transformation} $\gamma \colon (F,\eta_F) \Rightarrow (G, \eta_G) \colon M \to N$ of lax morphisms consists of $2$-cells $\gamma_c \colon F_c \Rightarrow G_c$ in $\K$ such that, writing $\gamma_{\bar{c}}\coloneqq \prod_i \gamma_{c_i}$, the diagram
\[
\begin{tikzcd}[ampersand replacement=\&]
N(B) \circ F_{\bar{c}} \arrow[d, "N(B)*\gamma_{\bar{c}}" description, Rightarrow] \arrow[r, "\eta^B_F", Rightarrow] \& F_d \circ M(B) \arrow[d, "\gamma_d*M(B)", Rightarrow] \\
N(B) \circ G_{\bar{c}} \arrow[r, "\eta_G^B"', Rightarrow]                                                           \& G_d \circ M(B)
\end{tikzcd}
\]
commutes for each functor symbol $B\colon \bar{c}\to d$ in $\Sigma$. Similarly, one obtains the $2$-category of $\Sigma$-models and oplax morphisms.

\begin{example}\label{ex:modificationsAndTransformations}
Let $F, G \colon \A \to \K$ be strict $2$-functors. We define the categories
\[
[F, G]_\mathbf{Lax}, \quad [F, G]_\mathbf{Oplax}, \quad [F, G]_\Psd, \quad [F, G]_\mathbf{Strict}
\]
as the categories whose objects are lax, oplax, strong, and strict morphisms $M_F \to M_G$ of $\Sigma_\A$-models in $\K$ (Example \ref{ex:2FunctorInucingAModel}), respectively, and whose morphisms are transformations between them. (The implicit size issues are sidestepped with Grothendieck universes.)

These objects correspond precisely to lax, oplax, pseudo-, and strict natural transformations $F \Rightarrow G$, and the morphisms between them are modifications.
\end{example}

\begin{definition}[The Object of Operations $M^\sigma$]
   Let $(\Sigma,\sigma)$ be a signature pair with set of sorts $S$, and let $M$ be a $\Sigma$-model in a $2$-category $\K$ with small products and inserters. We define $M_0 \coloneqq \prod_{s \in S} M_{c_s}$ and the \textbf{object $M^\sigma$ of operations} as the inserter
   \[
   \begin{tikzcd}
M^\sigma \arrow[d, "u_M"'] \arrow[r, "u_M"]                                    & M_0 \arrow[d, "(M(t_2))_f"]            \\
M_0 \arrow[r, "(M(t_1))_f"'] \arrow[ru, "\phi_M" description, Rightarrow] & \prod_{f} M_{c_f}
\end{tikzcd}
   \]
   where the indexing is over $\sigma$-function symbols $f\colon t_1\to t_2\colon c$.\footnote{For a term $t\in\Term_S^\Sigma$, we define $M(t)\colon M_0\to M_c$ recursively: $M(t) = \pi_s\colon M_0\to M_{c_s}$ if $t = s$ is a sort and $M(t) = M(B)(M(t_1),\ldots, M(t_n))$ if $t = B(t_1,\ldots, t_n)$.} For a term $t\colon c\in \Term_S^\Sigma$, we write $M^\sigma(t) \coloneqq M(t)\circ u_M\colon M^\sigma\to M_c$. If $\K = \CAT$, we call an object of $M^\sigma$ a \textbf{$\sigma$-model}.
\end{definition}
Explicitly, a $\sigma$-model $m$ in $M$ consists of an object $m_s$ of $M_{c_s}$ for each sort $s$ and a morphism $m(f)\colon M(t_1)((m_s)_s)\to M(t_2)((m_s)_s)$ for each function symbol $f\colon t_1\to t_2$.

\begin{definition}[Path Evaluation and the Object of Bialgebras $M^T$]\label{def:objectOfTBialgebras}
   Let $(\Sigma,\sigma)$ be a signature pair with a $\Sigma$-model $M$ in a PIE-complete $2$-category $\K$. We define a $2$-cell $\overline{p}_M \colon M^\sigma(t_1) \Rightarrow M^\sigma(t_2) \colon M^\sigma \to M_c$ recursively for each arrow $p \colon t_1 \to t_2 \colon c$ in the term graph $G_\sigma$:
   {\footnotesize
   \[
   \overline{p}_M = \begin{cases}
       \pi_f * \phi_M \colon M^\sigma(t_1) \Rightarrow M^\sigma(t_2), & \text{if $p = f \colon t_1 \to t_2$ is a function symbol in $\sigma$,} \\
       M(\alpha) * (M^\sigma(t_1'), \dots, M^{\sigma}(t_n')), & \text{if $p = \alpha_{t_1', \dots, t_n'}$ for a $\Sigma$-transformation symbol $\alpha^{\bar{x}}\colon t_1\to t_2$,} \\
       M(B) * (M^\sigma(t_1'), \dots, \bar{q}_M, \dots, M^\sigma(t_n')), & \text{if $p = B(t_1', \dots, q, \dots, t_n')$ for a $\Sigma$-functor symbol $B$.}
   \end{cases}
   \]
   }
   This extends to all paths $p$ in $G_\sigma$ by vertical composition; the empty path at a term $t$ is evaluated to the identity $2$-cell on $M^\sigma(t)$. The upcoming Lemma \ref{lem:substitution_evaluated} justifies that the definition of $\overline{p}_M$ is properly given by showing that the domains and codomains in the second recursive step agree, allowing the horizontal composition. 


   For a bialgebraic $\sigma$-theory $T$, the \textbf{object of $T$-bialgebras} $M^T$ is defined via the equifier
   \[
       \begin{tikzcd}[ampersand replacement=\&]
{M^T} \& {M^\sigma} \&\& {\prod_{(p,q)} M_c}
\arrow["{e_M}", from=1-1, to=1-2]
\arrow[""{name=0, anchor=center, inner sep=0}, "{(M^\sigma(t_1))}", shift left=3, from=1-2, to=1-4]
\arrow[""{name=1, anchor=center, inner sep=0}, "{(M^\sigma(t_2))}"', shift right=3, from=1-2, to=1-4]
\arrow["{(\bar{p}_M)}"', shift right=3, Rightarrow, from=0, to=1]
\arrow["{(\bar{q}_M)}", shift left=3, Rightarrow, from=0, to=1]
\end{tikzcd}
   \]
   where the indexing is over $(p,q\colon t_1\to t_2\colon c)\in T$. If $\K = \CAT$, we call an object of $M^T$ a \textbf{$T$-bialgebra}.
\end{definition}

\begin{proposition}\label{prp:bialgebraic_datum}
    Let $M$ be a $\Sigma$-model in a PIE-complete $\K$. Then $M^\sigma$ is the object of pre-bialgebras and $M^T$ the object of bialgebras of $\B_{M,T}$ over $M_0$. The four support morphisms of $\B_{M,T} = (M_0,\ f,g,\ f',g',\ \alpha,\beta)$ are obtained from morphisms $M(t)\colon M_0\to M_c$, where $t\colon c$ runs over the function arity, function coarity, equational arity, and equational coarity terms, respectively, and 
    $$
    \alpha = (\overline{p}_M)_{(p,q)}, \beta = (\overline{q}_M)_{(p,q)}\colon (M^\sigma(t_1))_{(p,q)}\Rightarrow (M^\sigma(t_2))_{(p,q)}\colon M^\sigma\to \prod_{(p,q)}M_c,
    $$
    where the indexing runs over $(p,q\colon t_1\to t_2\colon c)\in T$. 
\end{proposition}
\begin{proof}
    This is just a retelling of the definition.
\end{proof}

\begin{definition}\label{def:supportsOfTerms}
We call the morphism $M(t)\colon M_0\to M_c$ a \textbf{function/equational arity/coarity support} if $t\colon c\in \Term_S^\Sigma$ is a function/equational arity/coarity term, respectively, relative to $T$.
\end{definition}

\begin{remark}\label{rem:UPEM}
    By the $2$-categorical Yoneda lemma, $M^T$ is characterized up to isomorphism by an isomorphism of categories $\K(x,M^T)\cong\K(x,M)^T$, natural in $x$, where $\K(x,M)$ is the $\Sigma$-model in $\CAT$ induced by $M$. For example, let $(t\colon c\to c,m\colon tt\Rightarrow t,e\colon \id_c\Rightarrow t)$ be a monad in $\K$. Its (strict) Eilenberg--Moore object \cite{STREET1972149} is a triple $(d,u\colon d\to c,\eta\colon tu\Rightarrow u)$ satisfying $\eta\bullet mu=\eta\bullet t\eta$ and $\eta\bullet eu=\id_u$, such that:
    \begin{enumerate}
        \item for every triple $(x,h'\colon x\to c,\psi\colon th'\Rightarrow h')$ satisfying the same two equations, there is a unique $h\colon x\to d$ with $h' = uh$ and $\psi = \eta h$;
        \item for all $h_1,h_2\colon x\to d$ and every $\beta\colon uh_1\Rightarrow uh_2$ with $\eta h_2\bullet t\beta=\beta\bullet\eta h_1$, there is a unique $\theta\colon h_1\Rightarrow h_2$ with $u\theta=\beta$.
    \end{enumerate}
    The universal property of a general $M^T$ admits an analogous explicit description.
\end{remark}



\begin{lemma}\label{lem:substitution_evaluated}
   Let $(F,\eta) \colon M \to N$ be a lax morphism of $\Sigma$-models in $\K$. Let us denote by $V$ and $S$ the sets of variable symbols of $\Sigma$ and sorts of $\sigma$, respectively. Let $t\in \Term_V^\Sigma$ be a term of type $c$ in a context $\bar{x} = x_1\cdots x_n$, where $x_i\colon c_i$ for $i\leq n$. Let $t_1\colon c_1,\ldots, t_n\colon c_n\in \Term_S^\Sigma$. Then we have the following factorizations:
   \begin{align}
           M(t^{\bar{x}}(t_1,\ldots, t_n)) &\colon M_0\xrightarrow{(M(t_1),\ldots, M(t_n))}M_{\bar{x}}\xrightarrow{M_{\bar{x}}(t)} M_c\label{eq:0}\\
           \eta_{t^{\bar{x}}(t_1,\ldots, t_n)} &\colon N(t^{\bar{x}}(t_1,\ldots,t_n))F_0\xRightarrow{ N_{\bar{x}}(t)(\eta_{t_1},\ldots, \eta_{t_n})}N_{\bar{x}}(t)F_{\bar{x}}(M(t_1),\ldots, M(t_n))
           \nonumber\\
           &\quad \xRightarrow{\eta_{t,\bar{x}}(M(t_1),\ldots, M(t_n))}F_cM(t^{\bar{x}}(t_1,\ldots, t_n))\label{eq:1}
   \end{align}
\end{lemma}
\begin{proof}
We prove the claim by induction on the structure of the term $t$.

\noindent\textbf{Case 1:} Assume $t = x_i$. Then $t^{\bar{x}}(t_1,\ldots, t_n) = t_i$, $M(t_i) = \pi_i(M(t_1),\ldots, M(t_n))$ and $\eta_{t_i} = \id\bullet \pi_i(\eta_{t_1},\ldots, \eta_{t_n})$. Thus the equations hold.

\noindent\textbf{Case 2:} Assume $t = B(t_1',\ldots, t_k')$, where the equations hold for $t_i'$ for $i\leq k$ and $B$ is a functor symbol in $\Sigma$.
The substitution on $t$ then yields:
\[
t^{\bar{x}}(t_1, \dots, t_n) = B(t_1'^{\bar{x}}(t_1,\ldots, t_n),\ldots, t_k'^{\bar{x}}(t_1,\ldots, t_n))= B(s_1, \dots, s_k),
\]
where $s_i = t_i'^{\bar{x}}(t_1,\ldots, t_n)$ for $i\leq k$. By the induction hypothesis $M(s_i) = M_{\bar{x}}(t_i')\bar{M}$, where $\bar{M}\coloneqq (M(t_1),\ldots, M(t_n))$ for $i\leq k$.
To verify Equation~\eqref{eq:0}, consider:
\begin{align*}
   &M(t^{\bar{x}}(t_1, \dots, t_n)) \\
   &= M(B(s_1, \dots, s_k)) \\
   &= M(B)(M(s_1), \dots, M(s_k)) \tag{Definition of $M(B(s_1,\dots, s_k))$} \\
   &= M(B) (M_{\bar{x}}(t'_1)\bar{M}, \dots, M_{\bar{x}}(t'_k)\bar{M}) \tag{Induction hypothesis \eqref{eq:0} on $t'_i$} \\
   &= M(B) (M_{\bar{x}}(t'_1), \dots, M_{\bar{x}}(t'_k)) \bar{M} \\
   &= M_{\bar{x}}(B(t'_1, \dots, t'_k))\bar{M} \\
   &= M_{\bar{x}}(t)(M(t_1), \dots, M(t_n))
\end{align*}
This proves Equation~\eqref{eq:0} in the inductive step. Denote $\bar{\eta}\coloneqq (\eta_{t_1},\ldots, \eta_{t_n})$. For Equation~\eqref{eq:1}, consider:
{
\begin{align*}
   &\eta_{t^{\bar{x}}(t_1,\ldots, t_n)}\\
   &= \eta_{B(s_1,\ldots, s_k)}\\
   &= \eta^B (M(s_1),\ldots, M(s_k)) \bullet N(B)(\eta_{s_1},\ldots, \eta_{s_k}) \tag{Definition of $\eta_{B(s_1,\ldots, s_k)}$}\\
   &= \eta^B(M_{\bar{x}}(t_1')\bar{M},\ldots, M_{\bar{x}}(t_k')\bar{M})\bullet\\
   & N(B)(\eta_{t_1',\bar{x}}\bar{M} \bullet N_{\bar{x}}(t_1')\bar{\eta},\ldots, \eta_{t_k',\bar{x}}\bar{M}\bullet N_{\bar{x}}(t_k')\bar{\eta}) \tag{Induction hypothesis on $t_i'$}\\
   &=\eta^B(M_{\bar{x}}(t_1')\bar{M},\ldots, M_{\bar{x}}(t_k')\bar{M})\bullet
   N(B)(\eta_{t_1',\bar{x}}\bar{M},\ldots, \eta_{t_k',\bar{x}}\bar{M})\bullet\\
   & N(B)(N_{\bar{x}}(t_1')\bar{\eta},\ldots, N_{\bar{x}}(t_k')\bar{\eta}) \tag{Composition componentwise}\\
   &= (\eta^B(M_{\bar{x}}(t_1'),\ldots, M_{\bar{x}}(t_k'))\bullet
   N(B)(\eta_{t_1',\bar{x}},\ldots, \eta_{t_k',\bar{x}}))\bar{M} \bullet
    N(B)(N_{\bar{x}}(t_1'),\ldots, N_{\bar{x}}(t_k'))\bar{\eta} \\
   &= \eta_{B(t_1',\ldots, t_k'),\bar{x}}\bar{M}\bullet N_{\bar{x}}(B(t_1',\ldots, t_k'))\bar{\eta}\tag{Definitions of $\eta_{B(t_1',\ldots, t_k')}$ and $N_{\bar{x}}(B(t_1',\ldots, t_k'))$}\\
   &=\eta_{t,\bar{x}}(M(t_1),\ldots, M(t_n))\bullet N_{\bar{x}}(t)(\eta_{t_1},\ldots, \eta_{t_n})\tag{$t = B(t_1',\ldots, t_k')$}
\end{align*}
}
\end{proof}

\subsection{Examples}

\begin{examples}[Examples of $2$-Signatures]\label{ex:2-signatures}
   We define the following standard $2$-signatures representing background categorical structures.
   {\small
   \begin{align*}
       \Sigma_{\text{Cart}}
    &= (c,\quad \square\colon cc\to c,\ I\colon ()\to c,
    \quad x\xrightarrow{\delta^x}x\square x, \ x\xleftarrow{p_1^{xy}}x\square y\xrightarrow{p_2^{xy}}y,\ x\xrightarrow{!^x}I)\\
       \Sigma_{\text{Mon}}
       &= (c,\quad \square\colon cc\to c,\  I\colon ()\to c,\quad (x\square y)\square z\rightleftarrows x\square (y\square z),\ x\square I\rightleftarrows x\leftrightarrows I\square x)\\
       \Sigma_{\text{SMon}}
       &= (c,\quad \square\colon cc\to c,\ I\colon ()\to c,\quad (x\square y)\square z\rightleftarrows x\square (y\square z),\  x\square I\rightleftarrows x\leftrightarrows I\square x,\ x\square y\to y\square x)\\
       \Sigma_{\text{Adj}}
       &= (c,d,\quad L\colon c\to d,\ R\colon d\to c,\quad \eta^x\colon x\to R(L(x)),\ \varepsilon^y\colon L(R(y))\to y)\\
       \Sigma_{\text{Mnd}}
       &= (c,\quad T\colon c\to c,\quad \mu^x\colon T(T(x))\to T(x),\ \eta^x\colon x\to T(x))\\
       \Sigma_{\text{Distr}}
       &= (c,\quad T,S\colon c\to c,\quad  \mu_T^x,\ \eta_T^x,\ \mu_S^x,\ \eta_S^x,\ \gamma^x\colon S(T(x))\to T(S(x)))
   \end{align*}
   }
\end{examples}

\begin{examples}\label{exs:theories}
   We give the theories of monoids, Hopf algebras and fields, and a family of theories generalizing lax limits:
   \begin{enumerate}
       \item \textbf{Monoids:} Consider the signature pair $(\Sigma, \sigma)$ defined as follows:
       \begin{align*}
       \Sigma
       &= (c,\quad \square\colon cc\to c, I\colon ()\to c,\quad (x\square y)\square z\xrightarrow{\alpha^{xyz}} x\square (y\square z), x\square I\xrightarrow{\rho^x} x\xleftarrow{\lambda^x} I\square x)\\
       \sigma
       &= (s\colon c,\quad \mu\colon s\square s\to s, u\colon I\to s)
   \end{align*}
   For convenience, we write $x\square y$ for $\square(x,y)$. Consider the following parallel paths in the term graph $G_{\sigma}$ defining the theory of monoids $T$:
   \[
   \begin{tikzcd}[column sep=0.8 cm]
   (s \square s) \square s \arrow[d, "\mu \square s"'] \arrow[r, "\alpha_{s,s,s}"] & s \square (s \square s) \arrow[r, "s \square \mu"] & s \square s \arrow[d, "\mu"] & s \square I \arrow[r, "s \square u"] \arrow[rd, "\rho_s"'] & s \square s \arrow[d, "\mu"] & s \square s \arrow[d, "\mu"'] & I \square s \arrow[l, "u \square s"'] \arrow[ld, "\lambda_s"] \\
   s \square s \arrow[rr, "\mu"'] & & s & & s & s &
   \end{tikzcd}
   \]
   Let $M$ be a $\Sigma$-model in $\CAT$, for instance a monoidal category. We write $\otimes\coloneqq M(\square)$ and $I$ for $M(I)(*)$, where $*$ is the unique object of the terminal category. Note that
   \[
   M^\sigma = \Ins((M(s\square s),M(I)), (M(s),M(s))),
   \]
   where $M(s\square s)$, $M(s)$ and $M(I)$ are the functors
   \[
   x\mapsto x\otimes x, \quad x, \quad I\colon M_c\to M_c\text{, respectively.}
   \]
   Thus $M^\sigma$ consists of objects $(x,m\colon x\otimes x\to x, e\colon I\to x)$, and $M^\sigma$ can be understood as the category of pointed magmas of $M$. Let $U_M\colon M^\sigma\to M_c$ denote the forgetful functor. The internal monoids in $M$ form the full subcategory $M^T$ of $M^\sigma$.

   By definition, the category $M^T$ is the equifier of the parallel pair of natural transformations:
   \begin{align*}
   &\left(\overline{\mu \circ (\mu \square s)}_M, \,\, \overline{\rho_s}_M, \,\, \overline{\lambda_s}_M\right), \,\, \left(\overline{\mu \circ (s \square \mu) \circ \alpha_{s,s,s}}_M, \,\, \overline{\mu \circ (s \square u)}_M, \,\, \overline{\mu \circ (u \square s)}_M\right)\colon \\
   &\left(M((s \square s)\square s)U_M, \,\, M(s \square I)U_M, \,\, M(I \square s)U_M\right) \Rightarrow (U_M, \,\, U_M, \,\, U_M)\colon M^\sigma\to M_c^3
   \end{align*}
   This means an object $(x, m, e)$ of $M^\sigma$ is in $M^T$ if and only if
   {\small
   \[
   \left(\overline{\mu \circ (\mu \square s)}_M, \,\, \overline{\rho_s}_M, \,\, \overline{\lambda_s}_M\right)_{(x,m,e)} = \left(\overline{\mu \circ (s \square \mu) \circ \alpha_{s,s,s}}_M, \,\, \overline{\mu \circ (s \square u)}_M, \,\, \overline{\mu \circ (u \square s)}_M\right)_{(x,m,e)},
   \]
   that is, if and only if the following three equations hold in $M_c$:
   \begin{align*}
       m \circ (m \otimes x) &= m \circ (x \otimes m) \circ M(\alpha)_{x,x,x} \\
       m \circ (x \otimes e) &= M(\rho)_x \\
       m \circ (e \otimes x) &= M(\lambda)_x
   \end{align*}
   }
   These are precisely the classical associativity and unit laws for a monoid in $M$.

\item \textbf{Hopf Algebras:}
    To define internal Hopf algebras, we work relative to a skew-monoidal $2$-signature with a chosen interchange law:
    {\small
    \begin{align*}
        \Sigma&=(c, \quad \square \colon cc \to c,\ I \colon () \to c, \\
        &\quad (x \square y) \square z\xrightarrow{\alpha^{xyz}}  x \square (y \square z),\ I \square x \xrightarrow{\lambda^x} x \xrightarrow{\rho^x} x \square I,\ (x \square y) \square (z \square v) \xrightarrow{\chi^{xyzv}} (x \square z) \square (y \square v))
    \end{align*}
    }
    where $\chi^{xyzv}$ represents the interchange law, which is often uniquely constructed via the coherence of symmetric, or more generally, braided monoidal categories \cite{MacLane1998, joyalBraidedTensorCategories1993, streetQuantumGroupsPath2007}.

    We define the relative $1$-signature:
    \[
    \sigma_{\text{Hopf}} \coloneqq (s\colon c, \quad \mu\colon s \square s \to s, \ u\colon I \to s, \ \delta\colon s \to s \square s, \ \varepsilon\colon s \to I, \ S\colon s \to s)
    \]
    The theory $T_{\text{Hopf}}$ consists of the following equations in the term graph $G_{\sigma_{\text{Hopf}}}$, presented as commutative diagrams where each face represents a parallel pair of paths:\footnote{We use $p_1\square p_2$ for the path $(p_1\square t_2')\circ (t_1\square p_2)$, where $p_i\colon t_i\to t_i'$ is an arrow in the term graph $G_{\sigma_{\mathrm{Hopf}}}$ for $i = 1,2$.}
    {\small \[
\begin{tikzcd}[column sep = 0.7 cm]
(s\square s)\square s \arrow[d, "\mu\square s"'] \arrow[r, "{\alpha_{s,s,s}}"] & s\square (s\square s) \arrow[r, "s\square \mu"]              & s\square s \arrow[d, "\mu"]                                         & I\square s \arrow[rd, "\lambda_s"'] \arrow[r, "u\square s"]              & s\square s \arrow[d, "\mu" description]                                         & s\square I \arrow[l, "s\square u"']    &                        \\
s\square s \arrow[rr, "\mu"']                                                  &                                                              & s                                                                   &                                                                          & s                                                                               & s \arrow[l, "="'] \arrow[u, "\rho_s"'] &                        \\
s \arrow[d, "\delta"'] \arrow[r, "\delta"]                                     & s\square s \arrow[r, "\delta\square s"]                      & (s\square s)\square s \arrow[d, "{\alpha_{s,s,s}}"]                 & I\square s \arrow[d, "\lambda_s"']                                       & s\square s \arrow[l, "\varepsilon\square s"'] \arrow[r, "s\square \varepsilon"] & s\square I                             &                        \\
s\square s \arrow[rr, "s\square \delta"']                                      &                                                              & s\square(s\square s)                                                & s                                                                        & s \arrow[l, "="'] \arrow[u, "\delta" description] \arrow[ru, "\rho_s"']         &                                        &                        \\
(s\square s)\square (s\square s) \arrow[r, "{\chi_{s,s,s,s}}"]                 & (s\square s)\square (s\square s) \arrow[r, "\mu\square \mu"] & s\square s                                                          & I\square I \arrow[r, "\lambda_I"]                                        & I                                                                               & I\square I \arrow[r, "u\square u"]     & s\square s             \\
s\square s \arrow[u, "\delta\square \delta"] \arrow[rr, "\mu"']                &                                                              & s \arrow[u, "\delta"']                                              & s\square s \arrow[r, "\mu"'] \arrow[u, "\varepsilon\square \varepsilon"] & s \arrow[u, "\varepsilon"']                                                     & I \arrow[r, "u"'] \arrow[u, "\rho_I"]  & s \arrow[u, "\delta"'] \\
                                                                               & I                                                            & s\square s \arrow[rr, "s\square S"]                                 &                                                                          & s\square s \arrow[d, "\mu"]                                                     &                                        &                        \\
I \arrow[r, "u"'] \arrow[ru, "="]                                              & s \arrow[u, "\varepsilon"']                                  & s \arrow[r, "\varepsilon"] \arrow[u, "\delta"] \arrow[d, "\delta"'] & I \arrow[r, "u"]                                                         & s                                                                               &                                        &                        \\
                                                                               &                                                              & s\square s \arrow[rr, "S\square s"']                                &                                                                          & s\square s \arrow[u, "\mu"']                                                    &                                        &
\end{tikzcd}
\]
}

    Evaluating this theory in $\CAT$ over a symmetric monoidal category $M$ yields precisely the classical category $\Hopf_M$ of internal Hopf algebras.

\item \textbf{Fields:}
    The category of fields can be characterized starting from the category of commutative groups. Specifically, a field $k$ corresponds to a commutative group $k^*$ with a commutative ring structure on $k^*+1$ whose multiplication extends that of $k^*$.

    To formalize this, we consider a signature pair $(\Sigma, \sigma)$:
    \begin{align*}
    \Sigma
    &\coloneqq (c,\quad \square\colon cc\to c,\ I\colon ()\to c,\ S\colon c\to c,\\
    &\quad\quad x\xrightarrow{\delta^x}x\square x, \ x\xleftarrow{p_1^{xy}}x\square y\xrightarrow{p_2^{xy}}y,\ x\xrightarrow{!^x}I,\ x\xrightarrow{\eta^x}Sx)\\
    \sigma
    &\coloneqq (s \colon c,\quad m\colon s\square s\to s,\ i\colon s\to s,\ u\colon I\to s,\\
    &\quad\quad m_1,m_2\colon Ss\square Ss\to Ss,\ \nu\colon Ss\to Ss,\ \iota\colon I\to Ss)
    \end{align*}
    To define the bialgebraic theory of fields, we introduce the theory $T$ containing precisely the following equations in the term graph $G_\sigma$:
    \begin{itemize}
        \item The equations stating that $(s, m, i, u)$ satisfies the axioms of a commutative group.
        \item The equations stating that $(Ss, m_1, \nu, \iota, m_2, \eta_s \circ u)$ determines a commutative ring, where $(Ss, m_1, \nu, \iota)$ and $(Ss, m_2, \eta_s \circ u)$ satisfy the additive and the multiplicative equations, respectively.
        \item An equation stating that $\eta_s \colon s \to Ss$ determines a morphism of magmas from $(s,m)$ to $(Ss,m_2)$:
        \[
        m_2 \circ (\eta_s \square \eta_s) = \eta_s \circ m \colon s \square s \to Ss
        \]
    \end{itemize}
    We choose the categorical $\Sigma$-model $M = (\Set, \times, 1, (-)+1)$ where $M(\eta)$ is the unit of the maybe monad $(-)+1$. Then $M^T$ is equivalent to the category of fields.\footnote{One could take a bifunctor symbol $\boxplus$ with the codiagonal and injections instead of $S, \eta$ and $\iota$. This would be more advantageous in the sense that the coproduct can be defined as the left adjoint to the diagonal $1$-cell. This would allow an essentially unique formulation for fields if we were to develop the notion of $2$-equational $\Sigma$-theories.}

\item \textbf{Generalizing Lax Limits:}
    Consider an arbitrary $2$-signature $\Sigma = (C,\mathcal{F},\mathcal{T})$. We construct an associated algebraic theory $T$ for $\Sigma$ which generalizes the notion of a lax limit of a strict $2$-functor \cite{streetLimitsIndexedCategoryvalued1976, lack2CategoriesCompanion2010} to the cartesian (multivariable) setting. Specifically, this construction unifies the theories of monoids, commutative monoids, and monad algebras.

    We define the $1$-signature $\sigma_{\text{Lax}}$ relative to $\Sigma$:
    \[
    \sigma_{\text{Lax}} \coloneqq \left(s_c\colon c, \quad f_B\colon B(s_{c_1},\dots, s_{c_n})\to s_d\right)_{c\in C,B\colon c_1\cdots c_n\to d\in\mathcal{F}}.
    \]
    That is, the sorts are indexed by $C$ and the function symbols by $\mathcal{F}$. We recursively define for each term $t \colon c \in \Term_S^\Sigma$ a path $f_t \colon t \to s_c$ in the term graph $G_{\sigma_{\text{Lax}}}$:
    {\small
    \[
    f_t\coloneqq \begin{cases}
        ()_{s}\colon s\to s, & \text{if $t = s\in S$,}\\
        B(t_1, \dots, t_n) \xrightarrow{B(f_{t_1}, \dots, f_{t_n})} B(s_{c_1}, \dots, s_{c_n}) \xrightarrow{f_B} s_c, & \text{if $t = B(t_1,\dots, t_n)$, $t_i\colon c_i$ for $i\leq n$,}
    \end{cases}
    \]
    }
    where we use the convention that for any paths $p^i = (p^i_{m_i},\dots, p^i_1)\colon t_i'\to t_i''$ in $G_\sigma$ (for $i\leq n$) and a $\Sigma$-functor symbol $B$, the term $B(p^1, \dots, p^n)$ denotes the composition:
    \begin{align*}
        B(p^1,\dots, p^n) &\coloneqq B(p^1_{m_1},t_2'',\dots, t_n'')\circ\dots\circ B(p^1_1,t_2'',\dots, t_n'') \circ \cdots \\
        &\quad \dots \circ B(t_1',\dots, t_{n-1}',p^n_{m_n})\circ\dots\circ B(t_1',\dots, t_{n-1}', p^n_1).
    \end{align*}
    The theory $T_{\text{Lax}}$ consists of the following parallel paths:
    \[
    \begin{tikzcd}
    t_1^{\bar{x}}(\bar{s}) \arrow[r, "\alpha_{\bar{s}}"] \arrow[rd, "f_{t_1^{\bar{x}}(\bar{s})}"'] & t_2^{\bar{x}}(\bar{s}) \arrow[d, "f_{t_2^{\bar{x}}(\bar{s})}"] \\
                                                                                                   & s_c
    \end{tikzcd}
    \]
    for each $\alpha^{\bar{x}}\colon t_1\to t_2\in \mathcal{T}$, where $\bar{s} = s_{c_1}\cdots s_{c_n}$ and $x_i\colon c_i$ for $i\leq n$.

    If $\Sigma$ is the signature for a monad, a monoidal category or a symmetric monoidal category, then the associated algebraic theory is the theory of monad algebras, monoids and commutative monoids, respectively.
       \end{enumerate}

\end{examples}
\subsection{Induced Functor of Algebras}

In this subsection, we fix a PIE-complete $2$-category $\K$, a signature pair $(\Sigma,\sigma)$ and a bialgebraic $\sigma$-theory $T$. We find conditions under which a morphism $F\colon M\to N$ of $\Sigma$-models in $\K$ induces a morphism $F^T\colon M^T\to N^T$ between the objects of $T$-bialgebras. 

\begin{lemma}\label{lem:morphism_between_inserters}
   Let $(F,\eta) \colon M \to N$ be a lax morphism of $\Sigma$-models in $\K$. Assume that $\eta_t\colon N(t)F_0\Rightarrow F_c M(t)\colon M_0\to N_c$ is an isomorphism for each function coarity term $t\colon c$. Then the following hold:
   \begin{enumerate}
       \item There exists a unique morphism $F^\sigma\colon M^\sigma\to N^\sigma$ satisfying the equations
       \begin{align}
           u_N \circ F^\sigma
           &= F_0 \circ u_M \nonumber \\
           \phi_N * F^\sigma
           &= (\eta_{t_2}^{-1}*u_M)_f \bullet ((\prod_{f\colon t_1\to t_2\colon c} F_{c}) * \phi_M) \bullet (\eta_{t_1}*u_M)_f\label{eq:2}
   \end{align}
       \item For every path $p\colon t_1\to t_2$ in the term graph $G_\sigma$, the following diagram commutes:
       \[
\begin{tikzcd}
N(t_1)F_0u_M \arrow[r, "\eta_{t_1}*u_M", Rightarrow] \arrow[d, "\bar{p}_N*F^\sigma"', Rightarrow] & F_cM(t_1)u_M \arrow[d, "F_c*\bar{p}_M", Rightarrow] \\
N(t_2)F_0u_M \arrow[r, "\eta_{t_2}*u_M"', Rightarrow]                                             & F_cM(t_2)u_M
\end{tikzcd}
\]
In other words,
\begin{equation}\label{eq:3}
\eta_{t_2}u_M\bullet \bar{p}_NF^\sigma = F_c\bar{p}_M\bullet \eta_{t_1}u_M
\end{equation}

   \end{enumerate}
\end{lemma}
\begin{proof}
   The first universal property of the inserter $N^\sigma$ induces the unique morphism $F^\sigma$:
   \[
   \begin{tikzcd}[row sep = 1 cm]
M^\sigma \arrow[d, "u_M"'] \arrow[r, "u_M"]                                                 & M_0 \arrow[r, "F_0"] \arrow[d, "(M(t_2))_f" description]                                                                              & N_0 \arrow[dd, "(N(t_2))_f"] \\
M_0 \arrow[d, "F_0"'] \arrow[r, "(M(t_1))_f"'] \arrow[ru, "\phi_M" description, Rightarrow] & \prod_{f\colon t_1\to t_2\colon c}M_c \arrow[rd, "\prod_f F_c" description] \arrow[ru, "(\eta_{t_2}^{-1})_f" description, Rightarrow] &                              \\
N_0 \arrow[rr, "(N(t_1))_f"'] \arrow[ru, "(\eta_{t_1})_f" description, Rightarrow]          &                                                                                                                                       & \prod_f N_c
\end{tikzcd}
   \]
   We prove by structural induction on arrows $p$ in $G_\sigma$ that Equation~\eqref{eq:3} holds:

       \noindent\textbf{Case 1:} Assume $p = f$ is a $\sigma$-function symbol. Whiskering Equation~\eqref{eq:2} with $\pi_f$ from the left yields Equation~\eqref{eq:3}.

   \noindent\textbf{Case 2:} Assume $p = \alpha_{t_1',\ldots, t_n'}$, where $\alpha^{\bar{x}}\colon t_1''\to t_2''$ is a transformation symbol in $\Sigma$ and $t_i = t_i''^{\bar{x}}(t_1',\ldots, t_n')$ for $i = 1,2$.
   Now
     {\small \begin{align*}
       &\eta_{t_2}u_M\bullet \bar{p}_NF^\sigma\\
       &= \eta_{t_2}u_M\bullet N(\alpha)(N^\sigma(t_1'),\ldots, N^\sigma(t_n'))F^\sigma\tag{Definition of $\bar{p}_N $}\\
       &= \eta_{t_2}u_M\bullet N(\alpha)(N(t_1')F_0u_M, \ldots, N(t_n')F_0u_M)\tag{$N^\sigma(t_i')F^\sigma = N(t_i')u_N F^\sigma = N(t_i')F_0u_M$}\\
       &= (\eta_{t_2}\bullet N(\alpha)(N(t_1')F_0,\dots, N(t_n')F_0))u_M \tag{Whiskering}\\
       &= (\eta_{t_2'',\bar{x}}(M(t_1'),\ldots, M(t_n'))\bullet N_{\bar{x}}(t_2'')(\eta_{t_1'},\ldots, \eta_{t_n'})\bullet N(\alpha)(N(t_1')F_0,\dots, N(t_n')F_0))u_M \tag{\ref{eq:1}}\\
       &= (\eta_{t_2'',\bar{x}}(M(t_1'),\ldots, M(t_n'))\bullet N(\alpha)F_{\bar{x}}(M(t_1'),\dots, M(t_n'))\bullet N_{\bar{x}}(t_1'')(\eta_{t_1'},\ldots, \eta_{t_n'}))u_M\tag{Interchange}\\
       &= ((\eta_{t_2'',\bar{x}}\bullet N(\alpha)F_{\bar{x}})(M(t_1'),\ldots, M(t_n'))\bullet N_{\bar{x}}(t_1'')(\eta_{t_1'},\ldots, \eta_{t_n'}))u_M\tag{Whiskering}\\
       &= ((F_cM(\alpha)\bullet\eta_{t_1'',\bar{x}})(M(t_1'),\ldots, M(t_n'))\bullet N_{\bar{x}}(t_1'')(\eta_{t_1'},\ldots, \eta_{t_n'}))u_M\tag{Definition of lax morphism}\\
       &= (F_cM(\alpha)(M(t_1'),\ldots, M(t_n'))\bullet \eta_{t_1'',\bar{x}}(M(t_1'),\ldots, M(t_n'))\bullet N_{\bar{x}}(t_1'')(\eta_{t_1'},\ldots, \eta_{t_n'}))u_M \tag{Whiskering}\\
       &= (F_cM(\alpha)(M(t_1'),\ldots, M(t_n'))\bullet \eta_{t_1})u_M\tag{\ref{eq:1}}\\
       &= F_c\bar{p}_M \bullet \eta_{t_1}u_M\tag{Whiskering + Definition of $\bar{p}_M$}
   \end{align*}
   }

  \noindent \textbf{Case 3:} Assume $p = B(t_1',\ldots, q,\ldots, t_n')\colon B(t_1',\ldots, t_i',\ldots, t_n')\to B(t_1',\ldots, t_i'',\ldots, t_n')$, where $B\colon \bar{c} = c_1\cdots c_n\to c$ is a functor symbol in $\Sigma$ and the claim holds for the arrow $q\colon t_i'\to t_i''$. By definition, $\eta_{t_1}$ is the composite
  {\small\[
  N(B)(N(t_1'),\ldots, N(t_n'))F_0\xRightarrow{N(B)(\eta_{t_1'},\ldots, \eta_{t_n'})} N(B)F_{\bar{c}}(M(t_1'),\ldots, M(t_n'))\xRightarrow{\eta^B(M(t_1'),\ldots, M(t_n'))} F_cM(t_1)
  \]}
  Denote $M'' = (M^\sigma(t_1'),\ldots, M^\sigma(t_i''),\ldots,  M^\sigma(t_n'))$ and $M' = (M^\sigma(t_1'),\ldots, M^\sigma(t_n'))$. Now
  {\small\noindent
  \begin{align*}
       &\eta_{t_2}u_M\bullet \bar{p}_NF^\sigma\\
       &= \eta^B M''\bullet N(B)(\eta_{t_1'}u_M,\ldots,\eta_{t_i''}u_M,\ldots, \eta_{t_n'}u_M)\bullet\\
       & N(B)(N^\sigma(t_1')F^\sigma,\ldots,\bar{q}_N F^\sigma,\ldots, N^\sigma(t_n')F^\sigma) \tag{Definitions of $\eta_{t_2}$ and $\bar{p}_N$}\\
       &=  \eta^BM''\bullet N(B)(\eta_{t_1'}u_M, \ldots, \eta_{t_i''}u_M\bullet\bar{q}_NF^\sigma, \ldots, \eta_{t_n'}u_M)\\
       &= \eta^B M''\bullet N(B)(\eta_{t_1'}u_M,\ldots, F_{c_i}\bar{q}_M\bullet \eta_{t_i'}u_M,\ldots, \eta_{t_n'}u_M)\tag{Induction on $q$}\\
       &= \eta^B M''\bullet N(B)F_{\bar{c}}(M^\sigma(t_1'),\ldots,\bar{q}_M,\ldots, M^\sigma(t_n'))\bullet N(B)(\eta_{t_1'}u_M,\ldots, \eta_{t_n'}u_M)\\
       &= F_cM(B)(M^\sigma(t_1'),\ldots, \bar{q}_M,\ldots, M^\sigma(t_n'))\bullet \eta^B M'\bullet
       N(B)(\eta_{t_1'}u_M,\ldots, \eta_{t_n'}u_M)\tag{Interchange}\\
       &= F_c\bar{p}_M\bullet \eta_{t_1}u_M \tag{Definitions of $\bar{p}_M$ and $\eta_{t_1}$}
  \end{align*}
  }

  Thus Equation~\eqref{eq:3} holds for all arrows in $G_\sigma$, and it remains to prove it for paths. For the empty path it holds trivially. Let $p = p_n\cdots p_1$ be a path of arrows in $G_\sigma$. We treat the case $n = 2$; the general case follows by induction on $n$. Write $p\colon t_1'\xrightarrow{p_1}t_2'\xrightarrow{p_2}t_3'$. Then
  \begin{align*}
       &\eta_{t_3'}u_M\bullet \bar{p}_NF^\sigma\\
       &= \eta_{t_3'}u_M\bullet \bar{p_2}_NF^\sigma \bullet \bar{p_1}_NF^\sigma \tag{$p = p_2p_1$ + Whiskering}\\
       &= F_c\bar{p_2}_M\bullet \eta_{t_2'}u_M\bullet \bar{p_1}_NF^\sigma \tag{Equation~\eqref{eq:3} holds for $p_2$}\\
       &= F_c\bar{p_2}_M\bullet F_c\bar{p_1}_M\bullet \eta_{t_1'}u_M \tag{Equation~\eqref{eq:3} holds for $p_1$}\\
       &= F_c\bar{p}_M\bullet \eta_{t_1'}u_M \tag{$p = p_2p_1$ + Whiskering}
  \end{align*}\qedhere

\end{proof}
\begin{theorem}[Induced Functor of Algebras Theorem]\label{thm:induced_morphism_between_algebras}
   Let $(F,\eta)\colon M\to N$ be a lax morphism of $\Sigma$-models in $\K$. Assume the following:
\begin{enumerate}
   \item $\eta_t\colon N(t)F_0\Rightarrow F_cM(t)$ is an isomorphism for all function coarity terms $t\colon c$.
   \item $\eta_t*u_Me_M\colon N(t)F_0u_Me_M\Rightarrow F_cM(t)u_Me_M$ is a monomorphism in $\K(M^T,N_c)$ for every equational coarity term $t\colon c$.
\end{enumerate}
Then there is a unique morphism $F^T\colon M^T\to N^T$ such that $e_N F^T = F^\sigma e_M$.
\end{theorem}
\begin{proof}
   The first assumption with Lemma~\ref{lem:morphism_between_inserters}(1) yields the unique morphism $F^\sigma\colon M^\sigma\to N^\sigma$. Consider the following diagram, where the indexing is over $(p,q\colon t_1\to t_2\colon c)\in T$:
   \[
\begin{tikzcd}[ampersand replacement=\&]
 {M^T} \& {M^\sigma} \&\& {\prod_{(p,q)} M_c} \\
 \\
 {N^T} \& {N^\sigma} \&\& {\prod_{(p,q)} N_c}
 \arrow["{{{{e_M}}}}", from=1-1, to=1-2]
 \arrow["{{\exists! F^T}}"{description}, dashed, from=1-1, to=3-1]
 \arrow[""{name=0, anchor=center, inner sep=0}, "{{{{(M^\sigma(t_1))}}}}", shift left=3, from=1-2, to=1-4]
 \arrow[""{name=1, anchor=center, inner sep=0}, "{{{{(M^\sigma(t_2))}}}}"', shift right=3, from=1-2, to=1-4]
 \arrow["{{{F^\sigma}}}"', from=1-2, to=3-2]
 \arrow["{{{\prod_{(p,q)}F_c}}}", from=1-4, to=3-4]
 \arrow["{{e_N}}"', from=3-1, to=3-2]
 \arrow[""{name=2, anchor=center, inner sep=0}, "{{{{(N^\sigma(t_1))}}}}", shift left=3, from=3-2, to=3-4]
 \arrow[""{name=3, anchor=center, inner sep=0}, "{{{{(N^\sigma(t_2))}}}}"', shift right=3, from=3-2, to=3-4]
 \arrow["{{{{(\bar{p}_M)}}}}"', shift right=3, Rightarrow, from=0, to=1]
 \arrow["{{{{(\bar{q}_M)}}}}", shift left=3, Rightarrow, from=0, to=1]
 \arrow["{{{{(\bar{q}_N)}}}}", shift left=3, Rightarrow, from=2, to=3]
 \arrow["{{{{(\bar{p}_N)}}}}"', shift right=3, Rightarrow, from=2, to=3]
\end{tikzcd}
\]
   To induce the morphism $F^T$ using the equifier, we must show that $\bar{p}_N*F^\sigma e_M = \bar{q}_N*F^\sigma e_M$ for every $(p,q)\in T$. Let $(p,q\colon t_1\to t_2\colon c)\in T$. Then
   \begin{align*}
       \eta_{t_2}u_Me_M\bullet \bar{p}_NF^\sigma e_M
       &= F_c\bar{p}_Me_M\bullet \eta_{t_1}u_Me_M \tag{\eqref{eq:3}}\\
       &= F_c\bar{q}_Me_M\bullet \eta_{t_1}u_Me_M \tag{Equifier property of $e_M$}\\
       &= \eta_{t_2}u_Me_M\bullet \bar{q}_NF^\sigma e_M. \tag{\eqref{eq:3}}
   \end{align*}
   By the second assumption $\eta_{t_2}*u_Me_M$ is monic and hence $\bar{p}_N*F^\sigma e_M = \bar{q}_N*F^\sigma e_M$. Thus the equifier $N^T$ induces the unique morphism $F^T\colon M^T\to N^T$ with $e_NF^T = F^\sigma e_M$.
\end{proof}
\begin{corollary}
   Let $(F,\eta) \colon M \to N$ be a morphism of $\Sigma$-models in $\K$.
   \begin{enumerate}
       \item If $F$ is lax and $T$ is algebraic, then $F$ induces a $1$-cell $F^T \colon M^T \to N^T$ in $\K$.
       \item If $F$ is oplax and $T$ is coalgebraic, then $F$ induces a $1$-cell $F^T \colon M^T \to N^T$ in $\K$.
       \item If $F$ is strong, then $F$ induces a $1$-cell $F^T \colon M^T \to N^T$ in $\K$.
   \end{enumerate}
\end{corollary}
\begin{proof}
The lax case follows from Theorem~\ref{thm:induced_morphism_between_algebras}, since if $T$ is algebraic, then $\eta_t$ is an identity for each coarity term $t$ as coarity terms are sorts. The oplax case follows similarly. The strong case follows from Theorem~\ref{thm:induced_morphism_between_algebras}, since we inductively see that $\eta_t$ is an isomorphism for each term $t$.
\end{proof}

As the theory of monoids is algebraic, we recover the following classical result:
\begin{corollary}
   Let $F\colon M\to N$ be a lax monoidal functor between monoidal categories. Then there is an induced functor $\bar{F}\colon \Mon(M)\to \Mon(N)$ between the categories of internal monoids.
\end{corollary}
